\section*{Capítulo \ref{ch:g4:symmetry} --- Simetria}

\begin{solution}{exo:g4:symmetry:1}
Quadrado: sim ($4$ eixos). Retângulo: sim ($2$). Triângulo escaleno:
nenhuma dobra funciona. Círculo: sim --- qualquer reta que passe pelo
centro.
\end{solution}

\begin{solution}{exo:g4:symmetry:2}
Quadrado: $4$ (duas linhas médias e duas diagonais). Retângulo: $2$
(só as linhas médias). Triângulo equilátero: $3$ (uma por cada
\omterm{def:g4:geometry:polygon}{vértice}).
\end{solution}

\begin{solution}{exo:g4:symmetry:3}
Dobrando pela diagonal, as duas \omterm{def:g3:division:half}{metades} triangulares têm a mesma
forma, mas não caem uma sobre a outra --- uma delas está ``virada ao
contrário'': o decalque cai ao lado da figura, e não sobre ela. Logo
a diagonal não é um eixo.
\end{solution}

\begin{solution}{exo:g4:symmetry:4}
O L imagem fica a $2$ quadradinhos do \emph{outro} lado do eixo,
escrito ao contrário (um L espelhado), com os mesmos tamanhos.
\end{solution}

\begin{solution}{exo:g4:symmetry:5}
A figura completa mostra a bandeirinha e o reflexo dela de cabeça
para baixo abaixo do eixo, mastro debaixo de mastro.
\end{solution}

\begin{solution}{exo:g4:symmetry:6}
$P$ não se mexe: dobrar pelo eixo deixa todo ponto que está
\emph{sobre} o eixo exatamente onde estava --- $P$ é a imagem de si
mesmo.
\end{solution}

\begin{solution}{exo:g4:symmetry:7}
Eixo vertical: A, H, O, T. Eixo horizontal: B, E, H, O. Os dois: H,
O. Nenhum: F, N.
\end{solution}

\begin{solution}{exo:g4:symmetry:8}
Sim: o triângulo é isósceles, e o eixo dele passa pela ponta entre os
dois lados de $5$ cm e pelo meio da base de $3$ cm. A dobra troca os
dois lados iguais.
\end{solution}

\begin{solution}{exo:g4:symmetry:9}
Serve qualquer figura com exatamente dois eixos: um retângulo (não
quadrado), um sinal de mais com braços de dois comprimentos
diferentes (horizontal longo, vertical curto), uma elipse. Os dois
eixos sempre se cruzam no centro da figura.
\end{solution}

\begin{solution}{exo:g4:symmetry:10}
A casa completa é a \omterm{def:g3:division:half}{metade} esquerda mais a gêmea espelhada dela: uma
parede com o dobro da largura da meia parede, sob um telhado
triangular inteiro. Cada comprimento da direita é igual ao seu par da
esquerda --- a simetria copia comprimentos exatamente.
\end{solution}

\begin{solution}{exo:g4:symmetry:11}
Os desenhos dão razão nos dois sentidos: o triângulo isósceles dobra
pela reta que passa pela ponta (dois lados iguais $\to$ eixo); o
equilátero dobra de três maneiras; o escaleno não tem eixo (nenhum
lado igual $\to$ nenhum eixo). A dupla afirmação de Aya está certa
--- ela será demonstrada no \cref{ch:g6:symmetry}.
\end{solution}
